% CVT Foundations v0.2 — proposed Membrane Illustration replacement
% Review artifact. This section is intended to replace only the current
% membrane material and to stop immediately before the fusion section.

\section{Membrane Illustration: Regulated Transport Across an Interface}
\label{sec:membrane}

A regulated transport interface provides a bounded physical illustration of CVT. It is not presented as the universal mathematical form of the theory. The purpose of the illustration is narrower: to show how delivered flux, interface integrity, pre-existing capture readiness, productive uptake, restorative reserve, and delayed damage can be represented as distinct quantities in a system whose boundary is physically explicit.

\subsection{Established Interface-Transport Form}

Let domains \(A\) and \(B\) meet at an interface \(\Gamma\), and let \(c_A(x,t)\) and \(c_B(x,t)\) denote concentrations or other transport potentials on the two sides. Classical diffusion and membrane-transport models use boundary or interface laws that relate flux to a potential difference and an interface-transfer coefficient \cite{crank1975,kedem1958}. A linear permeability law may be written
\begin{equation}
    J_{\mathrm{raw}}(t)
    =
    \Pi_{\Gamma}(t)
    \bigl(c_B(t)-c_A(t)\bigr)
    \qquad \text{on }\Gamma,
    \label{eq:membrane-raw-flux}
\end{equation}
with the corresponding diffusive boundary flux on side \(A\), under a stated sign convention,
\begin{equation}
    -D_A\nabla c_A\cdot n_A
    =
    J_{\mathrm{raw}}.
    \label{eq:membrane-robin-law}
\end{equation}
Here \(D_A\) is a diffusion coefficient, \(n_A\) is the outward normal, and \(\Pi_{\Gamma}\) is an interface conductance or permeability with units chosen consistently with the transported quantity. Equations~\eqref{eq:membrane-raw-flux}--\eqref{eq:membrane-robin-law} are established transport forms, not CVT inventions. State dependence of \(\Pi_{\Gamma}\) is a modeling choice that must be justified by a physical mechanism or controller.

The interface law identifies what can cross. It does not by itself establish whether the receiving state can capture, retain, or survive the delivered input.

\subsection{Separating Permeability, Flux, Capture, and Uptake}

The membrane illustration distinguishes five quantities that should not be collapsed:
\begin{equation}
    \Pi_{\Gamma}
    \neq
    J_{\mathrm{raw}}
    \neq
    C
    \neq
    \rho_{\mathrm{obs}}
    \neq
    J_{\mathrm{eff}}.
    \label{eq:membrane-separation}
\end{equation}
The permeability \(\Pi_{\Gamma}\) is an interface property or control setting. The raw flux \(J_{\mathrm{raw}}\) is delivered transport. Capture readiness \(C\) is a pre-existing property of the receiving state, estimated before or independently of the uptake outcome. The quantity \(J_{\mathrm{eff}}\) is retained productive uptake measured over a declared accounting window. When raw and effective quantities are commensurate and delivered flux is positive, observed uptake efficiency is
\begin{equation}
    \rho_{\mathrm{obs}}
    =
    \frac{J_{\mathrm{eff}}}{J_{\mathrm{raw}}}.
    \label{eq:membrane-observed-uptake}
\end{equation}
This ratio is an outcome. It must not be used to define the same capture-readiness variable that is then claimed to explain it. In systems with endogenous amplification, storage delay, or noncommensurate input and output units, productive uptake requires an explicit conversion or target-response measure instead of Equation~\eqref{eq:membrane-observed-uptake}.

\subsection{Operational Mapping to \(B,Q,C,S\)}

For a declared receiving host, the membrane illustration maps the four candidate CVT conditions as follows:
\begin{enumerate}[label=(\roman*)]
    \item \textbf{bounded coupling \(B\)}: delivered flux remains within a two-sided admissible band relative to the host's changing handling capacity;
    \item \textbf{preserved identity-defining organization \(Q\)}: the interface and the host organization required by the success criterion remain functionally intact;
    \item \textbf{capture readiness \(C\)}: before the uptake outcome is observed, the receiving state lies in a region from which the specified input can be retained within an admissible basin over the declared horizon;
    \item \textbf{available restorative reserve \(S\)}: measurable recovery headroom remains after transport-induced strain.
\end{enumerate}

Let \(z_B,z_Q,z_C,z_S\) denote domain-native measurements and let \(G_i=g_i(z_i;\vartheta_i)\) be documented normalizations. The candidate membrane viability region is
\begin{equation}
    \mathcal K_{\mathrm{mem}}(t)
    =
    \mathcal K_B(t)
    \cap
    \mathcal K_Q(t)
    \cap
    \mathcal K_C(t)
    \cap
    \mathcal K_S(t).
    \label{eq:membrane-viability-region}
\end{equation}
This intersection is the representation-independent claim. When a scalar score is useful, candidate forms include
\begin{equation}
    \V_{\mathrm{prod,mem}}
    =
    G_BG_QG_CG_S,
    \qquad
    \V_{\min,\mathrm{mem}}
    =
    \min\{G_B,G_Q,G_C,G_S\},
    \label{eq:membrane-candidate-scores}
\end{equation}
and the wider aggregation family in Equation~\eqref{eq:aggregation-family}. No one form is privileged without comparative evidence against an independently declared viability outcome.

\subsection{Transport Regimes and Failure Signatures}

The same interface can enter several distinct regimes.

\paragraph{Insufficient delivery.}
Raw flux remains below the lower edge of the admissible operating band. The host may preserve organization and reserve while failing to receive enough input for the declared transformation. This is a bounded-coupling failure at the low-input side, not evidence that openness is always beneficial.

\paragraph{Overexposure.}
Raw flux exceeds the receiving system's handling band. Damage, rejection, spillover, saturation, or escape from the viable set may increase even when permeability and gradient remain physically large. Overexposure is first a failure of bounded coupling; observed uptake loss may follow but should be measured rather than assumed.

\paragraph{Capture failure.}
The interface delivers input while the receiving state lacks a viable capture path. A low value of \(
ho_{\mathrm{obs}}\) may be consistent with this mechanism, but it does not uniquely identify it: conversion losses, measurement error, transport delay, or an omitted state variable remain alternatives.

\paragraph{Loss of identity-defining organization.}
The interface or host organization required by the declared success criterion fails. Complete merger is a CVT failure only when continued distinction is part of that criterion. If merger is the intended transformation, the host and level of analysis must be redefined rather than treating all loss of boundary as pathological.

\paragraph{Reserve exhaustion.}
Transport remains possible and uptake may remain temporarily positive, but restorative reserve falls below the level required for repair or repeated-load survival. The failure may therefore appear only during a follow-up horizon after the target state has been reached.

These regimes show why permeability alone cannot classify membrane health. The same value of \(\Pi_{\Gamma}\) may be viable or destructive depending on the gradient, receiving state, damage, reserve, observation window, and declared success criterion.

\subsection{Causal State-Dependent Interface Control}

A regulated interface may adapt permeability to estimates of receiving-state condition. A causal policy class is
\begin{equation}
    \Pi_{\Gamma}(t)
    =
    \Pi_{\min}
    +
    \bigl(\Pi_{\max}-\Pi_{\min}\bigr)
    h_{\theta}\!\left(
        \widehat C(t-\Delta t),
        G_S(t-\Delta t),
        D(t-\Delta t),
        z_{\Gamma}(t-\Delta t)
    \right),
    \label{eq:membrane-causal-control}
\end{equation}
where \(0\leq h_{\theta}\leq1\), \(\Delta t\geq0\) is the sensing and actuation delay, and \(z_{\Gamma}\) contains declared interface measurements. A lagged value of \(
ho_{\mathrm{obs}}\) may enter \(z_{\Gamma}\) as feedback from a completed accounting window. Contemporaneous observed uptake cannot be treated as a pre-existing causal input without specifying how it is measured before the control action it is supposed to determine.

Equation~\eqref{eq:membrane-causal-control} is a controller class, not a universal membrane law. It should be compared with fixed permeability, capacity-maximizing control, and domain-native controllers using the same hosted-transformation outcome and follow-up horizon.

\subsection{Minimal Membrane Test}

A reproducible test should declare the host, transported quantity, interface law, viable set, transformation criterion, and follow-up horizon before fitting CVT variables. It should measure pre-input state and reserve, perturb permeability or driving gradient over a stated range, record delivered flux, retained productive uptake, rejection or damage, and recovery, and then compare product, minimum, soft-minimum, and direct-constraint representations out of sample. Deliberate counterexamples should include intended merger, successful transport with negligible measured reserve, high flux without uptake loss, and low capture estimates followed by stable productive retention.

The membrane illustration supports CVT only if the \(B,Q,C,S\) decomposition improves prediction or explanation beyond simpler transport, damage, or state-constraint models. A standard transport model that performs equally well with fewer independently measured quantities would constrain the added value of CVT in this domain.

\subsection{From the Interface Illustration to Fusion}

The next section uses artificial fusion as an illustrative boundary-mediated host problem. The transfer is structural: a plasma--machine system contains interfaces across which heat, particles, waves, fuel, ash, and control influence are transported or regulated while machine viability matters. This does not establish that fusion is literally reducible to Equation~\eqref{eq:membrane-robin-law}, that a tokamak is a biological membrane, or that the membrane illustration captures reactor physics. Those stronger claims require domain-native equations, calibrated parameters, and comparison with established fusion models.
